% ============================================================================
% CHAPITRE 4 — SECTION 4.1 : INTRODUCTION
% Fichier : partie4/introduction_partie3.tex

% ============================================================================

\section{INTRODUCTION}
\label{sec:introduction_resultats_discussions}
% ============================================================================
\begin{sloppypar}

Le cadre th\'eorique est maintenant pos\'e. Le chapitre~\ref{chap:contexte}
avait mis en \'evidence le paradoxe de l'authenticit\'e invisible,
formalis\'e six objectifs de recherche et construit l'outillage
th\'eorique n\'ecessaire. Les chapitres~\ref{chap:cadre_formel} et~\ref{chap:constructions} y avaient
r\'epondu par les \'etapes F, B, E, S et C de la progression
F$\to$B$\to$E$\to$S$\to$C$\to$X, en mobilisant le cadre CLO pour mesurer les compromis, le Trilemme CRO classique pour
en \'etablir les contraintes fondamentales, l'extension post-quantique
pour g\'en\'eraliser aux syst\`emes r\'esistants au quantum, la famille
SOTP comme ensemble de constructions de r\'ef\'erence dont l'optimalit\'e reste
conjecturale (Conjecture~\ref{conj:optimalite_universelle}) ainsi que la
formalisation Lean4 pour certifier le noyau alg\'ebrique. Le
chapitre~\ref{chap:experimentation} a sp\'ecifi\'e l'\'etape X.

La question qui se pose d\'esormais est diff\'erente. Les outils
th\'eoriques pr\'edisent-ils correctement le comportement des syst\`emes
r\'eels~? Les bornes d'impossibilit\'e correspondent-elles \`a des
limites observables~? Les hypoth\`eses institutionnelles sur lesquelles
repose le Trilemme, en particulier les valeurs de $\Delta$ et $\Delta_q$,
sont-elles coh\'erentes avec ce que l'on sait des capacit\'es cognitives
r\'eelles des institutions et des effets quantiques~? Ce chapitre
pr\'esente les \'el\'ements de r\'eponse disponibles, discute leurs
implications et trace les directions restant ouvertes.

Quatre axes structurent ce chapitre. Le premier est la calibration
illustrative du cadre sur des syst\`emes d\'eploy\'es, qui ne sont pas des
instances du th\'eor\`eme (Remarque~\ref{rem:portee_cas_etude}). Les ordres
de grandeur retenus y sont mutuellement coh\'erents, sans que cette
coh\'erence constitue une validation, tandis que les produits post-quantiques
illustratifs d\'epassent la borne. Le deuxi\`eme est la plausibilit\'e des
capacit\'es institutionnelles retenues, $\Delta \in [35, 73]$~bits et
$\Delta_q \in [25, 67]$~bits, au regard de la litt\'erature sur la
rationalit\'e born\'ee. Le troisi\`eme est le bilan de la formalisation
Lean4, qui comprend la compilation r\'eussie, l'hypoth\`ese de r\'egime \'etablie
comme n\'ecessaire, la comparaison chiffr\'ee entre la formulation initiale de
la borne et sa formulation certifi\'ee ainsi que la port\'ee exacte du t\'emoin de coh\'erence. Le quatri\`eme est l'expos\'e
des limitations, puisque la validation reste r\'etrospective et analytique en
attendant le programme exp\'erimental prospectif.

Deux types de validation doivent \^etre distingu\'es d\`es l'introduction. La coh\'erence interne, qui v\'erifie que les ordres de
grandeur retenus pour les syst\`emes existants ne se contredisent pas, est
ce que ce chapitre \'etablit, sans tester le cadre. La validation
exp\'erimentale prospective, qui mesurerait directement $\Delta$, $\alpha$,
$\beta$, $\Delta_q$ et $\eta_q$ aupr\`es d'acteurs institutionnels r\'eels
dans des conditions contr\^ol\'ees, constitue le programme de recherche
futur esquiss\'e \`a la section~\ref{sec:perspectives}.

La section~\ref{sec:resultats_theoriques} confronte les pr\'edictions du
cadre aux estimations analytiques et pr\'esente les r\'esultats formels de
la formalisation Lean4. La section~\ref{sec:discussions} interpr\`ete
ces r\'esultats, les positionne parmi les impossibilit\'es fondamentales
connues, examine leurs implications pour la conception et la
r\'eglementation, puis dresse le bilan des limitations et
probl\`emes ouverts. La section~\ref{sec:perspectives} trace enfin le
programme de recherche du sous-domaine naissant selon quatre axes
structur\'es.

\end{sloppypar}